%% ====================================================================
%%  Chapter 6 — Conclusion
%% ====================================================================
\chapter{Conclusion}
\label{ch:conclusion}

This thesis investigated whether the computational cost of DLoc was inherent to the task of localizing
devices via WiFi signals, or whether it represented over-parameterization of the model that could be
removed via model distillation or architectural alterations. The model was compressed into a family of
models of various sizes, trained both independently of DLoc and distilled from the DLoc model, and
evaluated along five dimensions: in-distribution accuracy, robustness, cross-session accuracy,
statistical validation, and computational cost. These evaluations lead to the following six findings.

First, it is possible to create a model that is extremely small relative to DLoc but still achieves the
same in-distribution accuracy; the MobileNetV2 model, with only 2.3~million parameters, achieves an
accuracy of 0.707\,m---a $7.2\times$ reduction in model size.

Second, the inclusion of skip connections in the MobileNetV2 model allows the resulting model to achieve
accuracy better than that of the DLoc baseline. An accuracy of 0.640\,m is achieved with the
MobileNetV2-UNet at 2.41~million parameters, a 9.5\% improvement over DLoc, with a model that is
$6.8\times$ smaller than DLoc.

Third, knowledge distillation acts as a method of increasing the robustness of the student model rather
than its accuracy. Accuracy is the same for the standalone and distilled models in the normal
in-distribution case, but the distilled models show improvements in the presence of noise ($3.4\times$
improvement) and in generalization to cross-session test folds.

Fourth, there is a minimum student-model size required to benefit from knowledge distillation. This
threshold is around 2~million parameters: the MobileNetV2-UNet (2.41M) and the MobileNetV2 (2.3M) lie
above it, but the TinyCNN, with only 47{,}000 parameters, lies below it, and distilling it degrades its
accuracy under AP dropout.

Fifth, the two methods of increasing the robustness of the model are complementary rather than additive.
Both the skip connections and knowledge distillation increase the robustness of the MobileNetV2-UNet, but
each helps most against a different type of corruption: the skip connections provide a $5.5\times$
improvement in noise robustness, while knowledge distillation most improves resilience under multi-AP
dropout. The two mechanisms are therefore complementary rather than additive.

Sixth, the cross-session accuracy of the MobileNetV2-UNet is competitive relative to the cost of the
model. Using the same leave-one-session-out protocol as DLoc, with three random seeds, the
MobileNetV2-UNet achieves median errors of 0.852\,m, 0.822\,m, and 1.116\,m on the three held-out August
sessions, averaging 0.930\,m against the published DLoc's 0.86\,m. It matches DLoc on the
different-furniture fold (0.822 vs.\ 0.82\,m) and remains level on the P90 metric, all at $6.8\times$
fewer parameters than DLoc.

These findings indicate that it is possible to achieve accurate WiFi-based localization by combining the
architectural design of the MobileNetV2-UNet (with skip connections) and the training strategy of
knowledge distillation. Each method resolves a different limitation: the skip connections resolve the
information loss that limits the spatial accuracy of the model, while knowledge distillation helps the
model generalize across distribution shifts in the WiFi signal. The distilled MobileNetV2-UNet achieves
an in-distribution accuracy of 0.640\,m, attains competitive cross-session accuracy, exhibits the best
results under multiple AP failures, and requires $56\times$ fewer floating-point operations and
$10\times$ less CPU inference time than DLoc. The cost of DLoc is therefore not inherent to WiFi-based
localization; a model at least an order of magnitude smaller than DLoc can meet or exceed its accuracy,
supporting its use in the edge devices for which WiFi-based indoor localization will be most useful.
